\section{Problem and data description}
\label{sec:data}

\subsection{Task and research question}
\reporttodo{Give the real-world motivation for image-text retrieval (for
example, photo search by description and description search by photo). Define
the two modalities: an RGB image $x\in\mathbb{R}^{3\times H\times W}$ and a
caption $t$ given as a token sequence. Define the prediction unit as one
image-caption pair, scored by a similarity $s(x,t)$. Image-to-text retrieval
ranks all test captions for each image, and text-to-image retrieval ranks all
test images for each caption. State the research question and the main
difficulties: many visually similar scenes, several valid descriptions per
image, and captions that also fit other images.}

\subsection{Dataset and study scope}
\reporttodo{Describe Flickr8k \cite{flickr8k}: source URL, version or access
date, licence status, collection and captioning process, number of images and
captions, the human judgement files, file formats, and storage size. Note the
dominance of dog and people scenes as a bias and limitation. Take the counts
from \codepath{results/a3/eda/flickr8k/summary.json}.}

\subsection{Dataset proposal and approval}
\label{sec:proposal}
\reporttodo{Summarize the Dataset Proposal submitted at Milestone 1
(\codepath{reports/a3_proposal.md}). Record the approval status (Approved,
Approved with conditions, or Rejected), the approval date, and any conditions.
State whether the retrieval threshold of about 5,000 image-text pairs
(Handbook Section 28.2) and the cross-modal validity constraints (Section 28.1)
are met.}

\subsection{Exploratory data analysis}
\label{sec:eda}
\reporttodo{Cover the multimodal EDA items in Handbook Section 31: pair counts,
missing modality rate, text length, and pair quality. Use
\codepath{results/a3/eda/flickr8k/summary.json} and explain what each figure
shows and why it matters for the models.}

% Fill from results/a3/eda/flickr8k/summary.json; do not type values by hand.
\begin{table}[htbp]
\centering\small
\caption{Flickr8k audit and EDA summary. Counts cover all images in the
release unless a split is named.}
\label{tab:eda-summary}
\begin{tabular}{lr}
\toprule
Check & Result\\
\midrule
Images / captions & \pending\\
Captions per image & \pending\\
Images that fail to decode & \pending\\
Images without captions / captions without an image & \pending\\
Captions with a malformed key & \pending\\
Images in no split file & \pending\\
Confirmed duplicate image pairs (crossing splits) & \pending\\
Training caption length: median / mean / max (tokens) & \pending\\
Training vocabulary / words seen once & \pending\\
Test tokens not seen in training captions & \pending\\
Test captions shared by several images & \pending\\
Test captions identical to a training caption & \pending\\
Own test pairs with a CrowdFlower majority yes & \pending\\
\bottomrule
\end{tabular}
\end{table}


\reportfigure{Images/flickr8k_pair_examples.png}{The first four images of the
official Flickr8k test list with their five reference captions.}{fig:pair-examples}

\reportfigure{Images/flickr8k_caption_length.png}{Caption length in tokens for
the train, dev and test splits, as a share of each split's
captions.}{fig:caption-length}

\reportfigure{Images/flickr8k_top_words.png}{The 20 most frequent content words
in training captions, with stop words removed.}{fig:top-words}

\reportfigure{Images/flickr8k_resolution_distribution.png}{Image sizes across
all 8,091 Flickr8k images. Marker area grows with the number of images of that
size.}{fig:resolutions}

\reporttodo{Describe pair quality: the CrowdFlower and expert judgements of
captions against their own images, captions shared by several images, and
captions written for another image that still fit. Explain how these cases make
Recall@K underestimate retrieval quality.}

\subsection{Splits and leakage prevention}
\label{sec:splits}
\reporttodo{State the split rule: the official Flickr8k train, dev and test
files, with dev used as validation and the test split held out. Give the split
unit (one image with all five of its captions), the images left out of every
split, and how leakage is prevented. State that the official split needs no
random seed or stratification, and give the seed used elsewhere, for example
for caption sampling and training. Cite the SHA-256 digests of the split files
(\texttt{split\_files} in \codepath{summary.json}). Report the duplicate check
across splits and the overlap that remains: dev and test captions whose text is
identical to a training caption (\texttt{identical\_to\_train\_caption}), and
different photographs of the same event. State the test isolation rule: models
and checkpoints are selected on dev, the test split is evaluated once after
configurations are fixed, and the human judgements are used only for
analysis.}

\begin{table}[htbp]
\centering\small
\caption{Flickr8k splits. The official split files are used unchanged; dev
serves as validation and the test split is held out.}
\label{tab:data-splits}
\begin{tabular}{lrrl}
\toprule
Split & Images & Image-caption pairs & Use\\
\midrule
Train & \pending & \pending & Optimization\\
Validation (dev) & \pending & \pending & Checkpoint selection\\
Test & \pending & \pending & Final evaluation only\\
\bottomrule
\end{tabular}
\end{table}


\reportfigure{Images/flickr8k_duplicate_examples.png}{The only confirmed
duplicate image pair. Both copies are in the training split.}{fig:duplicates}

\subsection{Preprocessing and augmentation}
\label{sec:preprocessing}
\reporttodo{Describe data cleaning and the handling of invalid annotations: the
captions whose key names no image, the images outside every split, and the
suspect pairs found in the EDA, with what is dropped and what is kept. Describe
image preprocessing (resize or crop to the encoder input
size, normalization constants, training augmentation) and text preprocessing
(tokenizer, lowercasing, maximum length, padding and truncation). Explain how
the five captions of an image are sampled during training. Give the image and
text tensor shapes and show one example batch after preprocessing.}

\reportfigure{Images/flickr8k_example_batch.png}{One training batch after
preprocessing: images with the caption sampled for each.}{fig:example-batch}
